% RESEARCH_DRAFT_SHA256: d1ffce57eeaa7185d20193500aba599b795886ebd5807bfa5a3178fe801e4648
% Layout and source-reference wrappers expanded for the web renderer.
\documentclass[11pt]{article}
\usepackage[margin=25mm]{geometry}
\usepackage{amsthm}
\usepackage{fontspec,unicode-math}
\setmainfont{Latin Modern Roman}
\setmathfont{Latin Modern Math}
\usepackage{amsmath,amssymb,mathtools}
\usepackage{xurl,hyperref}
\hypersetup{colorlinks=true,urlcolor=blue,linkcolor=blue}
\setcounter{secnumdepth}{0}
\setlength{\parindent}{0pt}
\setlength{\parskip}{5pt}
\setlength{\emergencystretch}{3em}
\newcommand{\sref}[2]{\hyperlink{source:#1:#2}{[S#2]}}
\newcommand{\cataloguescope}[1]{\paragraph{Source question.}#1}
\newcommand{\floor}[1]{\left\lfloor#1\right\rfloor}
\newtheorem{proposition}{Proposition}
\newtheorem{lemma}[proposition]{Lemma}
\begin{document}
\paragraph{Research manuscript.} OpenAI \texttt{gpt-6-astra}, reasoning \texttt{ultra}; 11 October 2026.
\paragraph{Preserved source body.} The complete catalogue mathematical text and sources follow with layout and source links expanded, including source qualifications and its historical preparation notes. New research begins at the section headed Research attempt.

% UnsolvedMath ID 30000607; source catalogue code OWR-1386-014
\section{Codegree Thresholds for Hypergraph Matchings}\label{30000607}
{\small UnsolvedMath ID 30000607 · Computer Science}
\subsection{Definitions and mathematical statement}

\paragraph{Shared notation.}Standard number systems are denoted by $\mathbb N,\mathbb Z,\mathbb Q,\mathbb R,\mathbb C$; index ranges, computational models and specialized notation follow the cited problem source and the definitions below.


A $k$-uniform hypergraph has edges that are $k$-element vertex subsets. A matching is a collection of pairwise disjoint edges. Its minimum $(k-1)$-codegree is the minimum number of edges containing a given $(k-1)$-subset.

A graph matching is a set of edges with pairwise disjoint endpoints, and it is perfect when every vertex is incident to one chosen edge. A weighted matching objective sums the chosen edge weights, unless an ordered objective is specified.
\paragraph{Statement recorded in the supplied source.}
For $n\geq k\geq2$, let $t_k(n)$ be the least $t$ such that every $n$-vertex $k$-uniform hypergraph $H$ with $\delta_{k-1}(H)\geq t$ has a matching of size $\lfloor n/k\rfloor$. If $k\nmid n$, is $t_k(n)=\lfloor n/k\rfloor$? \href{https://doi.org/10.4171/owr/2006/48}{[S2]}
\paragraph{Source status.}
The supplied catalogue labels this record open. Its source-status metadata was checked against the supplied snapshot on 2026-10-11. This import does not independently establish current open status or mathematical correctness. The catalogue literature assessment is dated 2026-08-21.
\subsection{Short English statement}
Determine the exact minimum-codegree threshold guaranteeing a largest possible matching when vertex count is not divisible by edge size.
\subsection{Sources}
\begin{itemize}
\item[S1] Ulam AI and the UnsolvedMath Contributors, UnsolvedMath: A Curated Collection of Open Mathematics Problems, record 30000607 (OWR-1386-014). Catalogue curation and metadata: CC BY 4.0. The supplied source question is retained; definitions and the English summary are editorial additions.
\url{https://huggingface.co/datasets/ulamai/UnsolvedMath}.
\item[S2] Combinatorics, Probability and Computing (2007), Oberwolfach Reports, DOI 10.4171/owr/2006/48, pp. 2875--2952
\url{https://doi.org/10.4171/owr/2006/48}.
\textit{Underlying Oberwolfach Reports material retains its source-specific licence, including CC BY-SA 4.0 where applicable; see THIRD\_PARTY\_NOTICES.md.}
\end{itemize}
\paragraph{Preparation and review.}
Original question attributed to the cited source and UnsolvedMath contributors. Codex (GPT-6; exact serving revision not exposed) assisted with definitions, English summaries, source inspection and formatting. No new human mathematical authorship or independent proof verification is asserted. Canonical ID 30000607 and source code OWR-1386-014 checked against the complete supplied registry and existing repository statements; known flagged repetitions and obvious identical conjectures excluded. Full mathematical-equivalence/literature review has not been performed. Statement preparation and final review followed contribution prompts C and D; automated checks verify structure and identity, not open status.
\label{end:30000607}

\section{Research attempt}
\subsection{Literature, source context and scope}
The full original catalogue body and source qualifications are preserved above. The relevant primary workshop passages were read in the cached publisher report. Historical claims and subsequent attributed results are distinguished from the mathematical work below. A restricted example, standard lemma or known theorem is not counted as new progress on the entire original target. Missing assumptions and repeated or rhetorical extracts are not silently turned into different solved problems.

\subsection{The sharp space barrier and the complete graph case}
Let $n\ge k\ge2$ with $k\nmid n$, and put $q=\lfloor n/k\rfloor$. The source conjectures minimum codegree $q$ suffices for a matching of size $q$. Later primary work by Han states this conclusion for each fixed $k\ge3$ and sufficiently large $n$; that external theorem is not reproved here and is not our new progress. We test the lower threshold and give a full proof of the elementary $k=2$ case.
\begin{proposition}
There is a $k$-uniform hypergraph with minimum codegree $q-1$ and largest matching size $q-1$. For $k=2$, every graph on $2q+1$ vertices with minimum degree at least $q$ has a matching of size $q$.
\end{proposition}
\begin{proof}
For the first assertion choose a set $A$ of $q-1$ vertices and take as edges all $k$-sets meeting $A$. A $(k-1)$-set disjoint from $A$ has exactly $q-1$ extensions to an edge. Such sets exist because the complement is large enough. A $(k-1)$-set meeting $A$ has every one of its $n-k+1$ possible extensions, and $n-k+1\ge q-1$. Thus the minimum codegree is $q-1$. Every matching uses a distinct vertex of $A$ per edge, so has size at most $q-1$. This size is attained by pairing each vertex of $A$ with its own $k-1$ distinct vertices of the complement; there are enough because $n-(q-1)\ge(k-1)(q-1)$. If $A$ is empty, the matching and edge set are both empty and the same conclusion holds.
For the graph assertion take a maximum matching of size $j$. If $j<q$, at least two vertices $u,v$ are unmatched. They have no unmatched neighbors, since an edge between unmatched vertices would enlarge the matching. For any matched edge $ab$, there are at most two edges from $\{u,v\}$ to $\{a,b\}$: three or four would include either $ua,vb$ or $ub,va$, permitting replacement of $ab$ by two disjoint edges. Hence $\deg u+\deg v\le2j<2q$, contradicting the minimum-degree assumption. Therefore $j=q$.
\end{proof}
\subsection{Verification and remaining gaps}
Every graph on three and five labeled vertices was checked for the stated degree implication, and the space-barrier hypergraph at $k=3,n=8$ was enumerated, giving minimum codegree 1 and matching number 1. The attempted extension of the graph augmentation proof to $k\ge3$ stops because neighbors of unmatched vertices no longer describe whole replacement hyperedges; intersections of larger candidate edges must also be controlled. No such general augmentation argument is supplied. The asymptotic higher-uniformity result remains attributed to Han, and no all-$n$ strengthening or independent new proof of the original conjecture is claimed.

\section{References and source attribution}
Primary workshop report: \url{https://doi.org/10.4171/owr/2006/48}. Relevant passages read: 2006/48 pp. 2928--2930, near-perfect matching threshold and Conjecture 4.

\url{https://arxiv.org/pdf/1404.1136}. Read Han’s primary Theorem 1.1 on near-perfect matchings: fixed k>=3, sufficiently large n not divisible by k, and minimum codegree at least floor(n/k). Its full proof is not independently reproduced here.

All retrievals were checked on 11 October 2026. Original question attribution remains in the preserved Sources section. External theorems are attributed and not claimed as new discoveries or fully verified merely by citation.


\paragraph{Integration review.} Codex (GPT-6; exact serving revision not exposed) reviewed the written arguments and expanded layout and source-reference wrappers. Repository metadata and the production web renderer were checked using Python 3.13 and JavaScript with MathJax 3.2.2. This model-assisted review is not independent human proof verification; no PDF compilation or formal proof checking is asserted.

\section{Completion Estimate}
\noindent\textbf{COMPLETION: 0\%}
No substantiated new progress toward the entire original target is claimed. The worked results and failed route are delimited above; a known result, illustrative model, or auxiliary obstruction is not counted as original-target completion.

\section{Final status and attribution}
\textbf{Attempt status: unresolved.} This describes this research attempt, not a new ruling on the source catalogue's status or on external literature claims.
The mathematics and this self-audit were generated by OpenAI \texttt{gpt-6-astra}, reasoning \texttt{ultra}, on 11 October 2026. Codex (GPT-6; exact serving revision not exposed) orchestrated the batch and contributed formatting/repository review. No human mathematical authorship, independent human proof checking, or formal verification is claimed. The historical preparation paragraphs in the reproduced catalogue body concern that body only.

Contribution prompt A was applied to the complete definition and supplied source record, followed by the adversarial checks of prompt B. Local reproduction files are \texttt{batch016/checks.py}, \texttt{computations.json}, and the per-ID audit JSON. They were not public links at preparation time. Numerical checks are finite and do not prove unrestricted results. No LaTeX engine was found locally; compilation was not performed by this research agent.

\end{document}
